\section{Harness Design for Agent-Assisted Formalization}\label{sec:workflow}

% The proposed harness-architecture figure is tracked in a separate figure
% issue.  Do not add a rendered figure or a figure reference until the
% author-supplied artwork is available.
%
% Historical workflow facts in this section are bound by CLM-007 to
% analysis/workflow-evolution.md and the six coded session indices in
% evidence/session-index/.  The architecture is an end-of-campaign synthesis,
% not a claim that every component was present throughout the campaign.

We reconstruct the campaign's end-state harness from controls introduced at
different stages. It separates scientific authority and implementation from
checks matched to the mathematical statement, Lean artifact, and public
declaration set, together with operational control and evidence capture.
% CLAIM: CLM-009
The resulting architecture is a synthesis of the final project state, not a
claim that every component operated throughout the campaign.

\subsection{Scientific authority and issue-scoped work}

The scientific owner retained authority over the mathematical scope. An issue
specified the intended result and acceptance conditions; for statements that
were easy to overgeneralize, a statement card recorded the hypotheses,
target conclusion, and permitted proof boundary before implementation. The
issue and statement card thus served different purposes: the issue bounded a
unit of work, while the card made the mathematical contract reviewable.
% CLAIM: CLM-009

The orchestrator coordinated work but did not edit Lean files. Implementing
agents received issue-scoped Git worktrees, separate working directories tied
to distinct branches, and returned changes through pull requests. The pull
request consolidated the committed candidate, validation results, and review
record. Coded records show that coordination and implementation were distinct
roles from the earliest observed session; later phases made the pull
request their durable boundary. % CLAIM: CLM-007, EV-2076, EV-0925

For the large refactoring campaign, a read-only scout examined repository
state and issue assumptions before an implementing agent was assigned.
Implementation and inspection therefore did not share edit authority. Two
review roles were also separated: one checked module structure, and another
compared public declarations before and after each refactor. This role
structure was observed in the July campaign; it should not be projected onto
earlier phases. % CLAIM: CLM-007, EV-1669

The minimal change contract was therefore as follows. The issue fixed scope
and acceptance conditions, a worktree held one implementation candidate, and
the pull request recorded the checks applicable to the changed artifact. The
scientific owner revised the statement contract after semantic findings;
implementers addressed code or build failures in the worktree. A candidate
advanced only after the applicable check results and a reviewer-attributed
completion record were attached to the pull request. % CLAIM: CLM-009

\subsection{Verification by artifact}

The harness assigned a distinct check to each artifact at risk.
% CLAIM: CLM-009

\paragraph{Mathematical statement.}
An agent independent of the implementer examined whether the proposition
matched the intended mathematics and whether its hypotheses were sufficient.
Reviewers tested candidate statements with concrete counterexamples when
appropriate. This semantic gate assessed the intended proposition and returned
scope corrections to the statement contract; the contrasting outcomes are
analyzed in Section~\ref{sec:incidents}.
% CLAIM: CLM-005, CLM-007, EV-0925

\paragraph{Lean artifact and axiom footprint.}
The build checked elaboration and type correctness. A separate
\texttt{\#print axioms} result recorded the assumptions of designated
declarations, while release-surface checks excluded experimental modules.
Together they certified the Lean artifact and its dependency boundary. A
release candidate then received a recorded full-build attestation for its
specific commit.
% CLAIM: CLM-009

\paragraph{Public declarations.}
Refactoring was checked against the prior public declaration set, independently
of downstream use. The declaration-level comparison certified interface
preservation by comparing the candidate inventory and, where exact
preservation was required, statement text. Section~\ref{sec:incidents}
describes the omission that motivated this boundary.
% CLAIM: CLM-007, EV-1669

\paragraph{Independent review record.}
A pull request required durable evidence that a reviewer distinct from the
implementer had examined the change. A dispatched review did not satisfy the
gate. The pull request needed an attributable record naming the reviewer and
result; the gate checked this completion evidence separately from the
substantive findings.
% CLAIM: CLM-007, EV-1669

\subsection{Ownership, liveness, and resource control}

Parallel work used independently named worktrees with one active implementing
owner per worktree. Shared import files were assigned to a single writer.
Silence alone did not transfer ownership: the orchestrator was expected to
check agent and process state before replacing an unresponsive worker. These
rules address two different risks, concurrent edits to shared state and
duplicated work caused by an incorrect liveness judgment.
% CLAIM: CLM-007, EV-2076, EV-1669, EV-0247

Lean builds shared a cache and ran under a global lock. Before a build, the
harness checked available memory across the container and removed stale helper
processes; under pressure it reduced Lean's thread count. These checks formed
the final operating arrangement and determined whether verification could
complete. Assurance about the mathematical statement came from the
checks above. % CLAIM: CLM-007, CLM-009, EV-0247

Agent assignments changed across the campaign. The durable role definitions
were therefore expressed through permissions and recorded outputs rather than
through model names. % CLAIM: CLM-007, EV-2076, EV-0925, EV-1669

\subsection{Evidence capture and observed events}

Session records, pull-request discussions, build results, and release
artifacts were retained as separate evidence sources. Claim identifiers link
paper statements to those sources, and reconstructed intervals are marked
separately from primary session evidence. This evidence layer supports audit
and later reconstruction alongside the verification artifacts.
% CLAIM: CLM-009

Across the campaign period from 2026-06-10 through the 2026-07-20 release,
the collected session logs record at least 15.1M output tokens, an
API-equivalent cost of \$6{,}973, and 199.1h of active engagement under a
5-minute inter-event gap cap. API-equivalent cost applies the published API
rates to recorded token classes; active engagement sums adjacent events only
when their separation is at most five minutes. The values are lower bounds
because the evidence store has known coverage gaps. % CLAIM: CLM-002

For the Vertex-hosted completion phase on July 2--3, 2026 (Pacific Time),
billing reconciliation kept three quantities separate. The provider export
measured a machine-wide actual charge of JPY~54{,}868; applying published
prices to the project logs yielded an unconditional calculation of
\(\$258.28\). Under the stated interpretation of the provider's aggregate
token metric, the project's derived share of the actual charge was
JPY~41{,}421--41{,}796. This attribution range is conditional on that metric
interpretation. % CLAIM: CLM-003

Six sessions were selected purposively for reconstruction value and coded by
the incident analyst against a fixed event vocabulary, separately from
evaluation of the formal result. Table~\ref{tab:event-codes}
describes only those sessions. The event families connect the evidence to the
design: statement and mathematical events call for semantic gates, whereas
ownership, handoff, and resource events call for operational and
review-completion controls.
% CLAIM: CLM-004

\begin{table}[t]
\centering
\caption{Observed events in the six purposively selected sessions.}\label{tab:event-codes}
\begin{tabular}{@{}p{0.72\linewidth}r@{}}
\toprule
Observed event & Count \\
\midrule
Statement mismatch or weakening & 4 \\
Recovery action & 19 \\
Resource failure & 12 \\
Worktree ownership conflict & 5 \\
Uncorrected false success report & 0 \\
\bottomrule
\end{tabular}
\end{table}

No uncorrected false success report, defined here as a mismatch
between a runtime success report and on-disk state, was observed in the coded
sessions. This is a descriptive count for the purposive sample.
% CLAIM: CLM-004
